\documentclass[12pt]{article}

%% Basic document formatting
\usepackage{amsmath, amsthm, amssymb, amsfonts}
\usepackage{mathtools}
\usepackage{xspace}
\usepackage{thmtools}
\usepackage{graphicx}
\usepackage{setspace}
\usepackage{fancyhdr}
\usepackage{titling}
\usepackage[left=0.4in,right=0.4in,top=1in,bottom=1in]{geometry}
\usepackage{float}
\usepackage{hyperref}
\usepackage{mathrsfs}
\usepackage{tabularx}
\usepackage[utf8]{inputenc}
\usepackage[english]{babel}
\usepackage{framed}
\usepackage[dvipsnames]{xcolor}
\usepackage{environ}
\usepackage{tcolorbox}
\tcbuselibrary{theorems,skins,breakable}

\usepackage{enumitem}
\setlist[enumerate]{leftmargin=*}
\setlist[enumerate,1]{labelindent=\parindent}
\setlist[enumerate,2]{labelindent=0pt}

% Blackboard Bold
\newcommand{\Z}{\mathbb{Z}}                 % integers
\newcommand{\Zpl}{\mathbb{Z}_{+}}           % positive integers
\newcommand{\Znn}{\mathbb{Z}_{\geq 0}}      % nonnegative integers. 
\newcommand{\Q}{\mathbb{Q}}                 % rationals
\newcommand{\Qpl}{\mathbb{Q}_{+}}           % positive Rationals
\newcommand{\R}{\mathbb{R}}                 % reals
\newcommand{\Rpl}{\mathbb{R}_{+}}           % positive Reals
\newcommand{\C}{\mathbb{C}}                 % complex numbers
\newcommand{\F}{\mathbb{F}}                 % field

% Words
\newcommand{\st}{\text{ such that }}
\newcommand{\wrt}{\text{ with respect to }}
\newcommand{\with}{\text{ with }}
\newcommand{\ie}{\text{, i.e. }}
\newcommand*{\ditto}{\text{---}\texttt{"}\text{---}}

% Operators
\newcommand{\abs}[1]{\left|#1\right|}                   % absolute value
\newcommand{\norm}[1]{\abs{\abs{#1}}}
\newcommand{\floor}[1]{\left\lfloor #1 \right\rfloor}   % floor
\newcommand{\ceil}[1]{\left\lceil #1 \right\rceil}      % ceiling
\newcommand{\powset}[1]{\mathscr{P}(#1)}

% Category Theory 
\renewcommand{\hom}{\text{Hom}}
\newcommand{\homspace}[3]{\hom_{#1}(#2, #3)}
\newcommand{\coproduct}[9]{
  \begin{tikzcd}[ampersand replacement=\&]
      \& \& #1 \arrow[dll, bend right, "#6"] \arrow[dl, "#5"] \\
   #4 \& #3 \arrow[l, "#9"] \\
      \& \& #2 \arrow[ull, bend left, "#8"] \arrow[ul, "#7"]
  \end{tikzcd}
}

\newcommand{\product}[9]{ 
  \begin{tikzcd}[ampersand replacement=\&]
      \& \& #3 \\
      #1 \arrow[r, "#7"] \arrow[urr, bend left, "#5"] \arrow[drr, bend right, "#6"] \& #2 \arrow[ur, "#8"] \arrow[dr, "#9"] \\ 
      \& \& #4
  \end{tikzcd}
}


% Algebra 
\newcommand{\Syl}[2]{\text{Syl}_{#1}{(#2)}}           % set of Sylow-p subgroups 
\newcommand{\<}{\langle}                % \<x\>, subgroup generated by x 
\renewcommand{\>}{\rangle}
\renewcommand{\index}[2]{[#1 : #2]}
\newcommand{\id}{\text{id}}             % identity element
\newcommand{\lhdneq}{%
  \mathrel{\ooalign{$\lneq$\cr\raise.22ex\hbox{$\lhd$}\cr}}}
\newcommand{\order}[1]{\text{o}(#1)}    % order of an element
\newcommand{\spec}{\operatorname{Spec}}
\newcommand{\rad}{\operatorname{rad}}   % radical of an ideal 
\newcommand{\sym}[1]{\mathfrak{S}_{#1}}
\newcommand{\aut}[1]{\text{Aut}(#1)} 
\newcommand{\gal}[2]{\text{Gal}(#1 / #2)} 
\newcommand{\inn}[1]{\text{Inn}(#1)} 
\let\oldcong\cong
\let\oldequiv\equiv
\renewcommand{\cong}{\oldequiv}
\renewcommand{\equiv}{\oldcong}

\newcommand{\triangleleftneq}{\mathrel{\ooalign{%
  $\trianglelefteq$\cr
  \hidewidth\raise0.06ex\hbox{$\not\mkern4mu$}\hidewidth\cr
}}}

\makeatletter % cycle
\newcommand{\cyc}[1]{(\mathbf{\cyc@process#1\relax})}
\def\cyc@process#1#2\relax{%
	#1%
	\ifx\relax#2\relax
	\else
		\,\cyc@process#2\relax
	\fi
}
\makeatother

\newcommand{\GL}[2]{\text{GL}_{#1}(#2)}                             % general linear group
\newcommand{\SL}[2]{\text{SL}_{#1}(#2)}                             % special linear group
\newcommand{\gl}[2]{\mathfrak{gl}_{#1}(#2)}
\renewcommand{\sl}[2]{\mathfrak{sl}_{#1}(#2)}
\newcommand{\so}[2]{\mathfrak{so}_{#1}(#2)}
\newcommand{\su}[1]{\mathfrak{su}(#1)}

% Linear Algebra
\newcommand{\bmat}[1]{\begin{bmatrix}#1\end{bmatrix}}   % bracketed matrix
\newcommand{\rank}{\operatorname{rank}}                 % rank
\newcommand{\nullity}{\operatorname{nullity}}           % nullity
\newcommand{\tr}{\operatorname{tr}}

% Combinatorics
\newcommand{\qnum}[1]{\left[#1\right]_q}                % q-analog
\newcommand{\qfact}[1]{\left[#1\right]!_q}              % q-analog factorial 
\newcommand{\qbinom}[2]{\binom{#1}{#2}_q}               % Gaussian binomial coefficient

% Topology/Analysis
\newcommand{\ball}[2]{\text{B}_{#1}(#2)}  % B_r(x): open r-balls around x
\newcommand{\diam}{\text{diam}}           % diamter of a set in metric space 
\newcommand{\lowint}[2]{\underline{\int_{#1}^{#2}}}
\newcommand{\upint}[2]{\overline{\int}_{#1}^{#2}}

% Misc Notation
\newcommand{\oneton}[1]{\{1, \ldots, #1\}}
\newcommand{\defeq}{\vcentcolon=}   % :=
\newcommand{\eqdef}{=\vcentcolon}   % =: 
\renewcommand{\bf}[1]{\textbf{#1}}
\newcommand{\im}{\operatorname{im}}
\newcommand{\fnrest}[2]{\left.#1\right|_{#2}}
\newcommand{\isom}{\overset{\sim}{\rightarrow}} 


\renewcommand{\thesection}{\arabic{section}.}
\renewcommand{\thesubsection}{(\alph{subsection})}

\newtheorem{claim}{Claim}
\newtheorem{theorem}{Theorem}
\newtheorem{proposition}{Proposition}
\newtheorem*{lemma}{Lemma}

%% Headers & title setup
\newcommand{\course}{Math140C}
\newcommand{\myname}{Jay Ser}
\setlength{\headheight}{14.5pt}
\pagestyle{fancy}
\fancyhf{}
\renewcommand{\headrulewidth}{0.4pt}
\lhead{\course}
\rhead{\myname}
\cfoot{\thepage}
\setlength{\droptitle}{-4em} 
\title{\course\ - HW \#5}
\author{\myname}
\date{2026.05.17} 

\begin{document}
\maketitle
\thispagestyle{fancy}

%------------------------------------------------------------------------------%
\newcommand{\M}{\mathcal{M}}
\newcommand{\scrC}{\mathscr{C}}
\newcommand{\sigalg}[1]{\sigma\text{-alg}(#1)}

\section{} 
\subsection{}
Briefly explain why every singleton set in $\R$ is Lebesgue measurable and has Lebesgue measure 0. 
Then, extend this result to every countable subset of $\R$. 

\noindent\dotfill 

Take any $x \in \R$. 
Letting 
$$A_n = (x - \frac{1}{n}, x + \frac{1}{n}), \quad (i \in \Zpl)$$
Recall $X \triangle Y$ for sets denotes the symmetric difference of $X$ and $Y$. 
For each $n$, 
\begin{equation} \label{eq:1_a_An}
  A_n \triangle \{x\} = (x - \frac{1}{n}, x) \sqcup (x, x + \frac{1}{n})
\end{equation}
The right hand side of \eqref{eq:1_a_An} is a finite union of intervals, so $A_n \triangle \{x\} \in \mathscr{E}$. 
Hence $m^{\star} (A_n \triangle \{x\}) = m(A_n \triangle \{x\})$. 
Additionally, it is clear that $m(A_n \triangle \{x\}) \rightarrow 0$ as $n \rightarrow \infty$, so
$$d(A_n, \{x\}) \rightarrow 0 \text{ as } n \rightarrow \infty$$
It immediately follows from this result that $\{x\}$ is Lebesgue measurable. 
In fact, more is true: 
For each $n \in \Zpl$, $A_n \in \mathscr{E}$ and is an open cover of $\{x\}$, which shows 
$$m^{*}(\{x\}) \leq m(A_n).$$
$m(A_n) \rightarrow 0$ as $n \rightarrow \infty$, so $m^{*} (\{x\}) = 0$, as was to be shown. 

The result for singletons can be used to show that every countable subset of $\R$ is Lebesgue measurable with Lebesgue measure 0:
if $U = \{x_1, x_2, \ldots\} \subset \R$ is any countable subset, $U$, being a countable (disjoint) union of singletons, is Lebesgue measurable because $\mathscr{M}(m)$ is a $\sigma$-algebra. 
Furthermore, since $m^{*}$ is countably additive on $\mathscr{M}(m)$, 
$$m^{*}(U) = \sum_{n \in \Zpl} m^{*} (\{x_n\}) = 0$$

\subsection{} 
Show directly from the definition of the Lebesgue measure that $m^{*}(\Q) = 0$. 

\noindent\dotfill 

Fix $\epsilon > 0$ and write $\Q = \{q_1, q_2, \ldots\}$. 
For each $n \in \Zpl$, let 
$$I_n = (q_n - \frac{\epsilon}{2^{n + 1}}, q_n + \frac{\epsilon}{2^{n + 1}})$$
Each $I_n$ has measure $m(I_n) = \frac{\epsilon}{2^n}$. 
$(I_n)_{n \in \Zpl}$ form an open cover for $\Q$, and 
$$\sum_{n = 1}^{\infty} m(I_n) = \sum_{n = 1}^{\infty} \frac{\epsilon}{2^{n}} = \epsilon,$$
as was to be shown. 

\subsection{} 
Show that $m^{*}(\R \setminus \Q) = \infty$ and $m^{*} (\Q \cap [0, 1]) = 1$. 

\noindent\dotfill

Since $\mathscr{M}(m)$ is a $\sigma$-algebra and $\Q$ is Lebesgue measurable, $\R \setminus \Q$ is also Lebesgue mesurable. 
Recall also that $\sigma$-algebras always contain the ambient set, so $\R$ is Lebesgue measurable. 
Because $m^{*}$ is monotonic on $\mathscr{M}(m)$ and there are bounded intervals of arbitrary length, it is clear that $m^{*}(\R) = \infty$. 
Using the countable additivity of $m^{*}$ on $\mathscr{M}(m)$, 
$$m^{*}(\R \setminus \Q) = m^{*} (\R) - m^{*} (\Q) = \infty,$$
as was to be shown. 

$[0, 1]$ and $\Q$ are Lebesgue measurable, so that $\mathscr{M}(m)$ is a $\sigma$-algebra means $[0, 1] \setminus \Q$ and $[0, 1] \cap \Q$ are also Lebesgue measurable. 
The quick computations 
\begin{gather*}
  m^{*}[0, 1] = m[0, 1] = 1 \\ 
  m^{*}(\Q \cap [0, 1]) = 0 \text{ (by monotonicity)} \\
  ([0, 1] \setminus \Q) \sqcup ([0, 1] \cap \Q) = [0, 1]
\end{gather*} 
give
$$m^{*} ([0, 1] \setminus \Q) = m^{*}[0, 1] - m^{*}(\Q) = 1.$$

\section{} 
\subsection{} 
Find a Lebesgue measurable set set $A$ such that $m(\overline{A}) > m(A)$. 

\noindent\dotfill 

Take $A = \Q$. 
$\overline{\Q} = \R$, and from Q1, 
$$m(\Q) = 0 < m(\overline{\Q}) = \infty$$

\subsection{} 
Find a Lebesgue measurable set $A$ such that $m(A^{\circ}) < m(A)$. 

\noindent \dotfill 

Let $A = \R \setminus \Q$, the set of irrational numbers.  
Since $\Q$ is measurable and measurable sets of a $\sigma$-algebra, $A$ is also measurable. 
Further, since $m$ is countably additive on measurable sets, by Q1, $m(A) = \infty$. 
$\Q$ is dense in $\R$ means $A^{\circ} = \emptyset$, and the fact $m(\emptyset) = 0$ means $A$ satisfies the desired properties. 

\section{} 
\subsection{} 
Prove that if $A, B \in \mathscr{E}$, then $\mu(B) = \mu(A \cap B) + \mu(B \setminus A)$. 

\noindent \dotfill 

As Brandon noted, if $A, B \in \mathscr{E}$, then $A \cap B \in \mathscr{E}$, which also means $B \setminus A \in \mathscr{E}$. 
$\mu$ is finitely additive on $\mathscr{E}$ and $B = (B \setminus A) \sqcup (A \cap B)$, giving 
$$\mu(B) = \mu(A \cap B) + \mu(B \setminus A)$$

\subsection{} 
Prove that if $B \subset \R^{p}$ and $A \in \mathscr{E}$, then $\mu^{*}(B) = \mu^{*} (B \setminus A) + \mu^{*} (B \cap A)$. 

\noindent \dotfill 

Let's first show $\mu^{*}(B \cap A) + \mu^{*} (B \setminus A) \leq \mu^{*}(B)$. 
Fix $\epsilon > 0$.
There is a collection $(A_n)_{n \in \Zpl}$ such that 
$$\text{each } A_n \in \mathscr{E} \text{ is open, } B \subset \bigcup_{n = 1}^{\infty} A_n, \text{ and } \sum_{n = 1}^{\infty} \mu(A_n) < \mu^{*}(B) + \epsilon$$
Since $\mu$ is regular, for each $n \in \Zpl$, pick $F_n, G_n \in \mathscr{E}$ such that $G_n$ contains $A$ and is open, $F_n$ is contained in $A$ and is closed, and 
$$\mu(A) - \epsilon / 2^n < \mu(F_n) \leq \mu(A) \leq \mu(G_n) < \mu(A) + \epsilon / 2^{n}.$$
Recall that intersection of elementary sets is also elementary. 
By construction of the collections $(E_n), \, (G_n)$, and $(F_n)$, 
\begin{gather*}
  \bigcup_{n = 1}^{\infty} (A_n \cap G_n) \text{ is an open cover for } B \cap A, \\  
  \bigcup_{n = 1}^{\infty} (A_n \setminus F_n) \text{ is an open cover for } B \setminus A
\end{gather*}
Also, for each $n \in \Zpl$, the following holds by the construction of $G_n$ and the monotocity of $\mu$: 
\begin{align*}
  \mu(A_n \cap G_n) & = \mu(A_n) + \mu(G_n) - \mu(A_n \cup G_n)  \\ 
					 & < \mu(A_n) + \mu(A) + \epsilon / 2^n - \mu(A_n \cup A) \\ 
					 & = \mu(A_n \cap A) + \epsilon / 2^n
\end{align*}
Hence 
\begin{align*}
  \mu^{*} (B \cap A) & \leq \sum_{n = 1}^{\infty} \mu(A_n \cap G_n)  \\ 
					 & < \sum_{n = 1}^{\infty} \mu(A_n \cap A) + \epsilon
\end{align*} 
Similarly, 
\begin{align*}
  \mu(A_n \setminus F_n) & = \mu(A_n \cup F_n) - \mu(F_n) \\ 
						 & < \mu(A_n \cup A) - \mu(A) + \epsilon / 2^n \\ 
						 & = \mu(A_n \setminus A) + \epsilon / 2^n,
\end{align*}
which gives 
$$\mu^{*}(B \setminus A) < \sum_{n = 1}^{\infty} \mu(A_n \setminus A) + \epsilon$$
Finally, since $(A_n \setminus A) \sqcup (A_n \cap A) = A_n$ for each $n$, 
\begin{align*}
  \mu^{*}(B \setminus A) + \mu^{*}(B \cap A) & < \sum_{n = 1}^{\infty} A_n + 2\epsilon \\ 
											 & < \mu^{*}(B) + 3\epsilon
\end{align*} 
Letting $\epsilon \rightarrow 0$, $\mu^{*} (B \setminus A) + \mu^{*} (B \cap A) \leq \mu^{*} (B)$.

Next, let's show $\mu^{*}(B) \leq \mu^{*} (B \cap A) + \mu^{*} (B \setminus A)$. 
Fix $\epsilon > 0$. 
There are collections $(A_n)_{n \in \Zpl}, (B_n)_{n \in \Zpl} \subset \mathscr{E}$ such that each $A_n$ and $B_n$ are open and 
\begin{gather*}
  B \setminus A \subset \bigcup_{n = 1}^{\infty} A_n, \quad \sum_{n = 1}^{\infty} \mu(A_n) < \mu^{*} (B \setminus A) + \epsilon \\ 
  B \cap A \subset \bigcup_{n = 1}^{\infty} B_n, \quad \sum_{n = 1}^{\infty} \mu(B_n) < \mu^{*} (B \cap A) + \epsilon
\end{gather*}
For each $n \in \Zpl$, since $\mu$ is regular, there are also $F_n, G_n \in \mathscr{E}$ such that $G_n$ contains $A$ and is open, $F_n$ is contained in $A$ and is closed, and 
$$\mu(A) - \epsilon / 2^{n + 1} < \mu(F_n) \leq \mu(A) \leq \mu(G_n) < \mu(A) + \epsilon / 2^{n + 1}.$$
For each $n \in \Zpl$, $A_n \setminus F_n$ and $B_n \cap G_n$ are still both elementary. 
Further, it is clear that $\bigcup_{n = 1}^{\infty} (A_n \setminus F_n) \text{ is an open cover for } B \setminus A$ and $\bigcup_{n = 1}^{\infty} (B_n \cap G_n) \text{ is an open cover for } B \cap A$, so  % FIXME: double check this? 
$$\bigcup_{n = 1}^{\infty} ((A_n \setminus F_n) \cup (B_n \cap G_n)) \text{ is an open cover for } B$$
For each $n \in \Zpl$, $(A_n \setminus F_n) \cap (B_n \cup G_n) \subset G_n \setminus F_n$, and $\mu(G_n \setminus F_n) = \mu(G_n) - \mu(F_n) < \epsilon / 2^n$. 
Using the monotocity of $\mu$, 
\begin{align*} 
  \mu((A_n \setminus F_n) \cup (B_n \cap G_n)) & \geq \mu(A_n \setminus F_n) + \mu(B_n \cap G_n) - \mu(G_n \setminus F_n) \\ 
											   & \geq \mu(A_n \setminus F_n) + \mu(B_n \cap G_n) - \epsilon / 2^{n} 
\end{align*} 
Putting all the information together and once again using the monoticity of $\mu$, 
\begin{align*}
  \mu^{*} (B) & \leq \sum_{n = 1}^{\infty} ((A_n \setminus F_n) \cup (B_n \cap G_n)) \\ 
			  & \sum_{n = 1}^{\infty} 
\end{align*} % FIXME: I would like both inequalities to flip

\subsection{} 

\section{} 
Let $X$ be a set and $\mathscr{M} \subset 2^{X}$ be a $\sigma$-algebra. 

\subsection{}
Take another set $Y$ and a function $h:X \rightarrow Y$. 
Show $\mathscr{M}' = \{B \subset Y \mid h^{-1}(B) \in \mathscr{M}\}$ is a $\sigma$-algebra on $Y$. 

\noindent\dotfill 

Take $B \in \mathscr{M}'$. 
Since $\mathscr{M}$ is a $\sigma$-algebra, $X \setminus h^{-1}(B) \in \mathscr{M}$. 
$$h^{-1}(Y \setminus B) = X \setminus h^{-1}(B),$$
so $Y \setminus B \in \mathscr{M}'$. 

Take $(B_n)_{n \in \Zpl} \subset \mathscr{M}'$. 
Then 
$$\bigcup_{n = 1}^{\infty} h^{-1}(B_n) = h^{-1}(\bigcup_{n = 1}^{\infty} B_n) \in \mathscr{M},$$
so $\bigcup_{n = 1}^{\infty} B_n \in \mathscr{M}'$. 
This shows $\mathscr{M}'$ is a $\sigma$-algebra on $Y$. 

\end{document}
